\paragraph{Cycle-depth wear sensitivity.}
As a final post hoc check, we rescored all 160 fixed matched-stock policy paths across the original and one-call contracts. We replaced the throughput and rate surcharges with whole-path rainflow damage using the NMC case-study curve $\Phi(d)=0.000524d^{2.03}$, normalized to 12 MWh nameplate, with both half and full cycles. Two cycle counters agreed on every depth histogram. Illustrative pack replacement prices of \$50, \$150 and \$300/kWh probe sensitivity to wear cost. No cell calibration or lifetime measurements are available.\footnote{B. Xu et al., \emph{Factoring the Cycle Aging Cost of Batteries Participating in Electricity Markets}, Eq.~24, \url{https://arxiv.org/abs/1707.04567}. Y. Shi et al., \emph{A Convex Cycle-based Degradation Model for Battery Energy Storage Planning and Operation}, Algorithm~1 and Table~I, \url{https://arxiv.org/abs/1703.07968}.}

At \$150/kWh, continuation saved \$50.26/\$31.10 per site-day against the price rule on the original seven-/thirty-day contract, while three-scenario dispatch was cheaper than continuation by \$7.32/\$3.90. On the one-call contract the corresponding continuation gains versus the rule were \$82.15/\$64.01; three-scenario dispatch saved \$3.07/\$2.54 versus continuation. The small seven-day one-call ranking reversed near \$59.16/kWh. Dispatch and services were fixed. New-optimum ceilings used the old lower bound less a rigorous maximum of old wear, plus the fixed incumbent's new cost, bounded using $d^{2.03}\le d^2$. These conservative thirty-day ceilings against continuation averaged \$687.24/day and \$662.16/day at \$150/kWh. No new rainflow-optimal dispatch or material-gain exclusion follows.
